% !PRINTED 19
\weberpage{19}

Ist dann $\Phi(x,y)$ eine ganze oder gebrochene rationale Funktion
von $x$ und $y$, und $y$ von $x$ durch die Gleichung (1) abhängig,
so sind
\weq{2}{\Phi(x,y)\,dx, \qquad \int \Phi(x,y)\,dx}
die zu der Kurve $F$ gehörigen \gesperrt{algebraischen Differentiale und
Integrale.}

Beispielsweise ist, wenn $f(x)$ eine Funktion dritten oder
vierten Grades von $x$ ist,
\[ y^2 - f(x) = 0 \]
die Gleichung einer Kurve dritten oder vierten Grades, zu der
die mit der Irrationalität $\sqrt{f(x)}$ behafteten \gesperrt{elliptischen
Differentiale und Integrale} gehören.

Wir wollen hier alle Kurven, deren Differentiale und Integrale
auf elliptische reduzierbar sind, \gesperrt{elliptische Kurven}
nennen. Wie das Beispiel zeigt, kommen darunter Kurven dritten
und vierten Grades vor.

Kegelschnitte gehören nicht zu den elliptischen Kurven, weil,
wenn zwischen $x$ und $y$ eine Gleichung zweiten Grades besteht,
$x$ und $y$ als rationale Funktionen eines Parameters $t$
dargestellt werden können, wodurch das algebraische Differential
auf ein rationales nach $t$ zurückgeführt werden kann. Solche
Kurven heißen \gesperrt{rationale Kurven}.

Wir werden sehen, daß alle Kurven dritten Grades ohne Doppel-
oder Rückkehrpunkt zu den elliptischen gehören. Kurven höheren
Grades können nur dann dazu gehören, wenn sie eine gewisse
Anzahl singulärer Punkte haben. Dies ist ein fundamentales
Kapitel in der allgemeinen Theorie der algebraischen Funktionen,
das nicht in den Plan dieses Werkes gehört\footnote{Die
wichtigsten diesen Gegenstand betreffenden Arbeiten sind:

\gesperrt{Riemann}, Theorie der \gesperrt{Abelschen} Funktionen [Crelles
Journal, Bd.\ 54 (1857)]. Mathematische Werke, 2.\ Aufl., S.\ 88,
487. Nachträge, herausgegeben von \gesperrt{Noether} und
\gesperrt{Wirtinger} (1902).

\gesperrt{Aronhold}, Monatsberichte der Berliner Akademie vom
25.\ April 1861.

\gesperrt{Clebsch}, Über die Anwendung der \gesperrt{Abelschen} Funktionen
in der Geometrie (Crelles Journal, Bd.\ 63).

\gesperrt{Clebsch}, Über diejenigen ebenen Kurven, deren
Koordinaten rationale Funktionen eines Parameters sind (ibid.,
Bd.\ 64).

\gesperrt{Clebsch}, Über diejenigen Kurven, deren Koordinaten
sich als elliptische Funktionen eines Parameters darstellen
lassen (ibid., Bd.\ 64).}.
